% =============================================================================
% ANNEXE J : Glossaire technique
% =============================================================================
% Consolide les définitions des termes techniques répétés dans le corps du
% portfolio (footnotes) afin d'alléger la lecture et de fournir une référence
% unique consultable de manière transversale.
% =============================================================================

\chapter{Glossaire technique}
\label{annex:glossaire}

\annexefiche
  {Consolidation produite en 2026 à partir des notes de bas de page du portfolio}
  {Référentiel transversal de vocabulaire technique utilisé dans les chapitres analytiques}
  {Lisibilité et précision terminologique (soutien transversal A2/B4/C2)}
  {Aucun caviardage (document pédagogique consolidé)}

Ce glossaire regroupe les définitions des termes techniques mobilisés dans le portfolio. Il vise à offrir une référence unique consultable de manière transversale, sans nécessiter le retour à chaque note de bas de page. Les termes sont classés par ordre alphabétique.

\bigskip

\begin{description}
  \item[\textbf{Agile (méthodes)}] Approches de gestion de projet fondées sur des cycles courts, la collaboration et l'adaptation continue. Elles privilégient des livraisons incrémentales plutôt qu'une planification figée.

  \item[\textbf{Architecture (logicielle)}] Organisation structurelle d'un système logiciel : découpage des composants, agencement, mise en relation, décisions techniques régissant leurs interactions (protocoles, interfaces, gestion des états).

  \item[\textbf{Artefact}] Document, outil ou livrable produit pendant l'action. Dans le cadre Tardif, sert de preuve observable de la compétence mobilisée.

  \item[\textbf{Backlog}] Liste priorisée des tâches, demandes et évolutions d'un projet. Mise à jour en continu selon les apprentissages et les contraintes.

  \item[\textbf{Board (agile)}] Représentation visuelle des tâches d'un projet, généralement structurée en colonnes (\textit{À faire / En cours / Terminé}). Donne une vue immédiate de l'avancement et des blocages.

  \item[\textbf{Branching}] Gestion des branches dans un système de versionnage : création de lignes de développement parallèles (typiquement une branche par fonctionnalité ou par version).

  \item[\textbf{Cascade (développement en)}] Approche séquentielle où chaque phase est clôturée avant la suivante. Risque principal : découvrir les erreurs tardivement, quand leur correction devient coûteuse. Synonyme : \textit{waterfall}.

  \item[\textbf{CI/CD}] \textit{Continuous Integration / Continuous Delivery (Deployment)}. Chaîne automatisée dans laquelle chaque modification de code déclenche la construction, les tests et, le cas échéant, le déploiement. Garantit que le code reste en permanence dans un état livrable et traçable.

  \item[\textbf{Commit}] Enregistrement d'un ensemble de modifications dans le système de versionnage, accompagné d'un message décrivant l'intention. Unité élémentaire de traçabilité du développement.

  \item[\textbf{Cycle de vie logiciel}] Ensemble des phases qu'un logiciel traverse, depuis l'expression du besoin jusqu'au retrait : analyse, conception, implémentation, vérification, validation, déploiement, maintenance.

  \item[\textbf{Definition of Done (DoD)}] Ensemble des critères qu'un livrable doit satisfaire pour être considéré comme terminé. En contexte médical, inclut le code, les tests, la documentation et la validation qualité.

  \item[\textbf{DevOps}] Ensemble de pratiques visant à rapprocher les équipes de développement (\textit{Dev}) et d'exploitation (\textit{Ops}), et à automatiser le cycle complet de livraison du logiciel.

  \item[\textbf{Exosquelette médical}] Dispositif mécatronique portable qui assiste ou remplace la motricité. En rééducation, guide les mouvements pour soutenir la reprise de la marche.

  \item[\textbf{Firmware}] Logiciel embarqué directement dans un composant matériel (microcontrôleur, carte électronique), au plus près du matériel. Pilote celui-ci en temps réel et n'est généralement pas modifiable par l'utilisateur final.

  \item[\textbf{Gate Review}] Revue de jalon. Point de contrôle formel à la fin d'une phase, durant lequel des critères prédéfinis sont vérifiés avant d'autoriser le passage à la phase suivante.

  \item[\textbf{GitLab}] Plateforme web intégrée de gestion de code source (basée sur Git) et d'industrialisation du développement logiciel : suivi des tickets, revues de code, intégration continue, registre d'artefacts.

  \item[\textbf{Hotfix}] Correction critique appliquée en urgence sur une version déjà livrée, en dehors du cycle normal de release.

  \item[\textbf{IEC 62304}] Norme internationale définissant les exigences relatives aux processus du cycle de vie des logiciels de dispositifs médicaux : développement, maintenance, gestion des risques, traçabilité. Obligatoire pour tout logiciel intégré dans un dispositif médical mis sur le marché européen.

  \item[\textbf{Ingénieur logiciel}] Professionnel qui conçoit, développe, teste, déploie et maintient des logiciels. À la différence du simple développeur, applique une démarche d'ingénieur : analyse formelle, choix raisonné d'architecture, prise en compte des contraintes de fiabilité, performance, sécurité et coût.

  \item[\textbf{Itératif}] Démarche procédant par cycles successifs produisant des livrables testables. Chaque cycle intègre des retours pour ajuster le suivant.

  \item[\textbf{JFrog Artifactory}] Gestionnaire d'artefacts logiciels : entrepôt centralisé qui stocke, versionne et trace les binaires produits par le développement (bibliothèques, paquets, images de conteneurs).

  \item[\textbf{KPI}] \textit{Key Performance Indicator}. Indicateur clé de performance, mesurant la pertinence d'un choix ou d'une activité selon des critères opérationnels explicites.

  \item[\textbf{Lean (méthode)}] Approche visant à maximiser la valeur produite tout en réduisant les gaspillages. En logiciel, privilégie des incréments utiles et rapides plutôt que des livraisons monolithiques.

  \item[\textbf{Livrable}] Tout produit, document ou artefact formellement remis à un destinataire (client, partie prenante, autorité réglementaire) au terme d'une activité ou d'un jalon.

  \item[\textbf{Logiciel embarqué}] Programme s'exécutant directement sur des composants matériels (microcontrôleurs, processeurs embarqués). Responsable du pilotage en temps réel, de l'acquisition des données capteurs et de la communication entre modules.

  \item[\textbf{MDR}] \textit{Medical Device Regulation} (Règlement UE 2017/745). Règlement européen sur les dispositifs médicaux, entré en application en mai 2021. Renforce les exigences en matière de sécurité, performance clinique, surveillance post-marché et traçabilité.

  \item[\textbf{Médiation}] Processus structuré où un tiers neutre aide deux parties à converger sans imposer de décision. Se distingue de l'arbitrage, où la décision est tranchée par une autorité.

  \item[\textbf{Pipeline (CI/CD)}] Chaîne d'étapes automatisées déclenchées à chaque modification du code source : compilation, tests, analyse de qualité, packaging, déploiement.

  \item[\textbf{Prototype}] Version partielle d'un produit destinée à tester une hypothèse de conception ou un besoin utilisateur. Niveau de fidélité variable selon l'objectif.

  \item[\textbf{Release}] Livraison officielle d'une version du logiciel, figée et traçable.

  \item[\textbf{Rétrospective}] Réunion de fin de sprint où l'équipe analyse ce qui a fonctionné, ce qui a bloqué, et les actions à mettre en place. Alimente l'amélioration continue.

  \item[\textbf{SAFe for Medical}] \textit{Scaled Agile Framework for Medical Devices}. Adaptation de SAFe aux contraintes réglementaires du médical, conciliant cadence agile et traçabilité normative.

  \item[\textbf{Scrum}] Cadre de travail (\textit{framework}) agile organisant le développement par sprints (2 à 4 semaines). Définit trois rôles (Product Owner, Scrum Master, Équipe), quatre événements (Planning, Daily, Review, Rétrospective), trois artefacts (Product Backlog, Sprint Backlog, Increment).

  \item[\textbf{SDP}] \textit{Software Development Plan}. Document de planification du cycle de vie logiciel obligatoire dans le contexte IEC~62304 (clause 5.1).

  \item[\textbf{SOUP}] \textit{Software of Unknown Provenance}. Composants logiciels tiers non développés en interne. En contexte médical, leur suivi (version, origine, risque, licence) est requis par IEC~62304.

  \item[\textbf{Sprint}] Période de travail courte et fixe (typiquement 2 semaines), au terme de laquelle l'équipe livre un incrément testable. Unité de base de Scrum.

  \item[\textbf{Sprint Planning}] Réunion d'ouverture du sprint durant laquelle l'équipe sélectionne les éléments du backlog et planifie le travail.

  \item[\textbf{Start / Stop / Continue}] Format de rétrospective en trois volets : ce qu'on commence, ce qu'on arrête, ce qu'on poursuit. Transforme le bilan en plan d'action concret.

  \item[\textbf{Template}] Gabarit. Modèle réutilisable qui standardise la structure d'un document. Son usage réduit le temps de production et les écarts de forme entre livrables.

  \item[\textbf{Timeboxing}] Attribution d'une durée fixe et non extensible à une activité. En facilitation, évite l'enlisement et force la production d'un résultat dans le temps imparti.

  \item[\textbf{V-Model}] Modèle de développement structuré en V. Branche descendante : conception (exigences, architecture, conception détaillée). Branche ascendante : validation par tests correspondants (unitaire, intégration, système, acceptance). Très utilisé en contexte réglementé.

  \item[\textbf{Vélocité}] En agile, quantité moyenne de travail (souvent mesurée en \textit{story points}) qu'une équipe parvient à terminer par sprint. Indicateur de capacité de production, non de valeur livrée.

  \item[\textbf{Versionnage}] Gestion systématique des différentes versions successives d'un fichier ou d'un logiciel. Permet de retrouver l'état exact à une date donnée, de comparer des versions, ou de revenir en arrière. Exigence fondamentale de traçabilité en contexte médical.

  \item[\textbf{Vote par points}] Méthode de sélection collective où chaque participant dispose d'un nombre limité de votes. La convergence s'appuie sur une agrégation des préférences du groupe.
\end{description}

